\subsection{EXPONENTIAL MAPPING}

THEOREM 4. Let \(\mathbf{G}\) be a Lie group. There exists one and only one exponential mapping of \(\mathbf{G}\) defined on \(\mathrm{L}(\mathbf{G})\) .

There exist a convex open neighbourhood \(\mathbf{U}\) of 0 in \(\mathrm{L}(\mathbf{G})\) and an exponential mapping of \(\mathbf{G}\) defined on \(\mathbf{U}\) . It can be assumed, by choosing \(\mathbf{U}\) sufficiently small, that

\[
\phi ((\lambda + \lambda^ {\prime}) a) = \phi (\lambda a) \phi (\lambda^ {\prime} a)\tag{1}
\]

for \(a \in \mathrm{L}(\mathbf{G})\) , \(\lambda, \lambda'\) in \(\mathbf{K}\) , \(\lambda a, \lambda'a\) , \((\lambda + \lambda')a\) in \(\mathbf{U}\) .

Let \(a \in \mathrm{L}(\mathrm{G})\) . There exists an integer \(n > 0\) such that \(\frac{1}{n} a \in \mathrm{U}\) . If \(m\) is another integer such that \(\frac{1}{m} a \in \mathrm{U}\) , then \(\frac{1}{nm} a \in \mathrm{U}\) and relation (1) implies

\[
\phi \left(\frac {1}{n} a\right) = \left(\phi \left(\frac {1}{n m} a\right)\right) ^ {m}, \quad \phi \left(\frac {1}{m} a\right) = \left(\phi \left(\frac {1}{n m} a\right)\right) ^ {n};
\]

hence \(\left(\phi\left(\frac{1}{n}a\right)\right)^n = \left(\phi\left(\frac{1}{m}a\right)\right)^m\) . There exists an extension \(\psi: \mathrm{L}(G) \to G\) of \(\phi\) such that \(\psi(a) = \left(\phi\left(\frac{1}{n}a\right)\right)^n\) for \(a \in \mathrm{L}(G)\) and \(n\) an integer \(>0\) such that \(\frac{1}{n}a \in \mathrm{U}\) . Clearly \(\psi\) is analytic and is an exponential mapping of \(G\) . If \(\psi': \mathrm{L}(G) \to G\) is an exponential mapping of \(G\) , \(\psi\) and \(\psi'\) coincide on a neighbourhood of 0 and hence are equal since \(\mathrm{L}(G)\) is connected.

Henceforth when we speak of the exponential mapping of G, we shall mean the mapping considered in Theorem 4. It will be denoted by \(\exp_{G}\) or exp if there is no risk of confusion.

Example. Let A be a complete normed unital associative algebra. Then \(\exp_{\mathbf{A}}\) is the exponential mapping defined in Chapter II, § 7, no. 3.

PROPOSITION 7. Let G be a Lie group and a an element of L(G). The mapping \(\lambda \mapsto \exp(\lambda a)\) of K into G is the unique morphism \(\phi\) of the Lie group K into G such that \((\mathrm{T}_0\phi)1 = a\) .

The mappings \((\lambda, \lambda') \mapsto \exp(\lambda a) \exp(\lambda' a)\) and \((\lambda, \lambda') \mapsto \exp(\lambda + \lambda') a\) of \(K \times K\) into \(G\) are analytic and coincide on a neighbourhood of \((0, 0)\) . As \(K \times K\) is connected, these mappings are equal. Hence \(\phi: \lambda \mapsto \exp(\lambda a)\) is a Lie group morphism of \(K\) into \(G\) . The tangent mapping at \(0\) to \(\lambda \mapsto \lambda a\) is the mapping \(\lambda \mapsto \lambda a\) ; and \(T_0(\exp) = \mathrm{Id}_{\mathrm{L}(\mathrm{G})}\) ; hence \((T_0\phi)1 = a\) . The uniqueness assertion of the proposition follows from Theorem 1.

PROPOSITION 8. Let G be a Lie group. For all x, y in L(G) and n an integer,

\begin{enumerate}
\def\labelenumi{(\arabic{enumi})}
\setcounter{enumi}{1}
\tightlist
\item
\end{enumerate}

\[
\exp (x + y) = \lim _ {n \rightarrow + \infty} \left(\left(\exp \frac {1}{n} x\right)\left(\exp \frac {1}{n} y\right)\right) ^ {n}\tag{2}
\]

\[
\exp [ x, y ] = \lim _ {n \rightarrow + \infty} \left(\left(\exp \frac {1}{n} x\right)\left(\exp \frac {1}{n} y\right)\left(\exp \frac {1}{n} x\right) ^ {- 1} \left(\exp \frac {1}{n} y\right) ^ {- 1}\right) ^ {n ^ {2}}\tag{3}.
\]

By Proposition 7, this follows from Proposition 4 of § 4, no. 3, taking \(\lambda = \frac{1}{n}\) .

PROPOSITION 9. Let G be a complex Lie group and G' the underlying real Lie group. Then \(\exp_{G} = \exp_{G'}\) .

This follows from Proposition 5 of §4, no. 3 and the analyticity of \(\exp_{G}\) and \(\exp_{G'}\) .

PROPOSITION 10. Let G and H be Lie groups and \(\phi\) a morphism of G into H.

\begin{enumerate}
\def\labelenumi{(\roman{enumi})}
\item
  \(\phi \circ \exp_{\mathrm{G}} = \exp_{\mathrm{H}} \circ \mathrm{L}(\phi)\) .
\item
  If G is an integral subgroup of H, then \(\exp_{G} = \exp_{H} | L(G)\) .
\end{enumerate}

The two sides of the equality (i) are analytic mappings of L(G) into H which coincide in a neighbourhood of 0 (§ 4, Proposition 8, no. 4) and hence are equal. Assertion (ii) is a special case of (i).

COROLLARY 1. Let G be a Lie group, G' a Lie subgroup of G and \(a \in \mathrm{L}(\mathrm{G})\) . The following conditions are equivalent:

\begin{enumerate}
\def\labelenumi{(\roman{enumi})}
\item
  \(a \in \mathbf{L}(\mathbf{G}')\) ;
\item
  \(\exp (\lambda a)\in \mathbf{G}'\) for all \(\lambda \in \mathbf{K}\) and \(|\lambda |\) sufficiently small;
\item
  \(\exp(\lambda a) \in \mathbf{G}'\) for all \(\lambda \in K\) .
\end{enumerate}

The argument is as in § 4, no. 4, Corollary 1 to Proposition 8.

COROLLARY 2. Let \(\mathbf{G}\) be a Lie group, \(\mathbf{H}\) an integral subgroup of \(\mathbf{G}\) and \(a \in \mathbf{L}(\mathbf{G})\) . Consider the following conditions:

\begin{enumerate}
\def\labelenumi{(\roman{enumi})}
\item
  \(a \in \mathrm{L}(\mathrm{H})\) ;
\item
  \(\exp_{\mathbf{G}}(\lambda a)\in \mathbf{H}\) for all \(\lambda \in \mathbf{K}\) .
\end{enumerate}

Then (i) \(\Rightarrow\) (ii). If the topology of H admits a countable base, then (i) \(\Leftrightarrow\) (ii).

Let \(i\) be the canonical injection of \(\mathbf{H}\) into \(\mathbf{G}\) . If \(a \in \mathbf{L}(\mathbf{H})\) , then

\[
\exp_ {\mathrm{G}} (\lambda a) = \left(\exp_ {\mathrm{G}} \circ \mathrm{L} (i)\right) (\lambda a) = (i \circ \exp_ {\mathrm{H}}) (\lambda a) \in \mathrm{H}.
\]

Hence (i) \(\Rightarrow\) (ii). The converse when H has a countable base follows from Proposition 3.

COROLLARY 3. Let G be a Lie group, \(\rho\) an analytic linear representation of G, \(x \in L(G)\) and \(g \in G\) .

\begin{enumerate}
\def\labelenumi{(\roman{enumi})}
\item
  \(\rho(\exp x) = \exp L(\rho)x;\)
\item
  \(\mathrm{Ad}(\exp x) = \exp \mathrm{ad} x;\)
\item
  \(g(\exp x)g^{-1} = \exp(\operatorname{Ad} g.x).\)
\end{enumerate}

The argument is as in § 4, no. 4, Corollaries 2 and 3 to Proposition 8.

COROLLARY 4. Let G be a finite-dimensional connected Lie group.

\begin{enumerate}
\def\labelenumi{(\roman{enumi})}
\item
  \(\text{Int}(L(G)) = \text{Ad}(G).\)
\item
  Let \(\mathbf{Z}\) be the centre of \(\mathbf{G}\) . Then \(\mathbf{Z}\) is a Lie subgroup of \(\mathbf{G}\) whose Lie algebra is the centre of \(\mathbf{L}(\mathbf{G})\) . The mapping of \(\mathbf{G} / \mathbf{Z}\) into \(\operatorname{Int} \mathbf{L}(\mathbf{G})\) derived from \(g \mapsto \operatorname{Ad} g\) when passing to the quotient is a Lie group isomorphism.
\end{enumerate}

Assertion (i) follows from Corollary 3 (ii) and the remarks after Definition 2. Let \(g \in G\) . In order that \(\operatorname{Ad} g = \operatorname{Id}_{\operatorname{L(G)}}\) , it is necessary and sufficient that \(\operatorname{Int} g\) coincide with \(\operatorname{Id}_G\) on a neighbourhood of \(e\) (§ 4, no. 1, Theorem 1 (ii)) and hence on the whole of \(G\) ; in other words, it is necessary and sufficient that \(g \in Z\) . Then (ii) follows from Corollary 1 to Proposition 1.

DEFINITION 3. Let G be a finite-dimensional connected Lie group. The Lie group \(\operatorname{Int}(\mathbf{L}(\mathbf{G})) = \operatorname{Ad}(\mathbf{G})\) is called the adjoint group of G.

PROPOSITION 11. Let \(\mathbf{G}\) be a connected commutative Lie group.

\begin{enumerate}
\def\labelenumi{(\roman{enumi})}
\item
  exp is an étale morphism of the additive Lie group L(G) onto G.
\item
  If \(\mathbf{K} = \mathbf{R}\) and \(\mathbf{G}\) is finite-dimensional, \(\mathbf{G}\) is isomorphic to a Lie group of the form \(\mathbf{R}^p \times \mathbf{T}^q\) ( \(p, q\) integers \(\geqslant 0\) ).
\end{enumerate}

By the Hausdorff formula, \((\exp x)(\exp y) = \exp (x + y)\) for \(x,y\) sufficiently close to 0 and hence for all \(x,y\) in \(\mathrm{L}(G)\) by analytic continuation. Hence exp is a group homomorphism and is étale since

\[
\mathrm{T} _ {e} (\exp) = \mathrm{Id} _ {\mathrm{L(G)}}.
\]

Hence (i). Assertion (ii) follows from (i) and General Topology, Chapter VII, § 1, Proposition 9.

PROPOSITION 12. Let G be a Lie group and L = L(G). For all \(x \in L\) , let \(\mathrm{T}_{x}(\mathrm{L})\) be identified with L, so that the right differential \(\varpi(x)\) of exp at x is a linear mapping of L into L. For all \(x \in L\) ,

\[
\varpi (x) = \sum_ {n \geqslant 0} \frac {1}{(n + 1) !} (\mathrm{ad} x) ^ {n}.
\]

The two sides are analytic functions of \(x\) and are equal for \(x\) sufficiently close to 0 (§ 4, no. 3, Proposition 6).

Remark. \(\varpi(x) . (\mathrm{ad} x) = \exp \mathrm{ad} x - 1\) . We write, by an abuse of notation,

\[
\varpi (x) = \frac {\exp \operatorname{ad} x - 1}{\operatorname{ad} x}.
\]

COROLLARY. Let \(\mathbf{G}\) be a complex Lie group and \(x \in \mathrm{L}(\mathbf{G})\) . The tangent mapping at \(x\) to \(\exp_{\mathbf{G}}\) has kernel \(\bigoplus_{n \in \mathbf{Z}\setminus\{0\}} \operatorname{Ker}(\operatorname{ad} x - 2i\pi n)\) .

The integral function \(z \mapsto \sum_{n \geqslant 0} \frac{1}{(n + 1)!} z^n\) , equal to \(\frac{e^z - 1}{z}\) for \(z \neq 0\) , admits

as zeros the points of \(2\pi i\mathbf{Z}\setminus\{0\}\) , which are all simple zeros. The corollary then follows from Proposition 12 and the following lemma:

Lemma 2. Let E be a complex Banach space, u an element of \(\mathcal{L}(\mathrm{E})\) , S the spectrum of \(u\) in \(\mathcal{L}(\mathrm{E})\) (Spectral Theories, Chapter I, §1, no. 2) and \(f\) a holomorphic complex function on an open neighbourhood \(\Omega\) of S. Suppose that \(f\) admits in \(\Omega\) only a finite number of distinct zeros \(z_{1},\ldots ,z_{n}\) , of multiplicities \(h_1,\dots,h_n\) . Then \(\operatorname {Ker}f(u)\) is the direct sum of the \(\operatorname {Ker}(u - z_i)^{h_i}\) for \(1\leqslant i\leqslant n\) .

(For the definition of \(f(u)\) , see Spectral Theories, Chapter I, § 4, no. 8.)

There exists a holomorphic function \(g\) on \(\Omega\) , everywhere non-zero, such that \(f(z) = (z - z_1)^{h_1} \ldots (z - z_n)^{h_n} g(z)\) . Then \(g(u)g^{-1}(u) = g^{-1}(u)g(u) = 1\) and hence \(\operatorname{Ker}f(u) = \operatorname{Ker}\prod_{i=1}^{n}(u - z_i)^{h_i}\) . Considering \(\operatorname{Ker}f(u)\) as a \(\mathbf{C}[\mathbf{X}]\) -module by means of the external law \((h,x)\mapsto h(u)x\) for \(h\in \mathbf{C}[\mathbf{X}]\) , \(x\in \operatorname{Ker}f(u)\) , we see that \(\operatorname{Ker} f(u)\) is the direct sum of the \(\operatorname{Ker}(u - z_i)^{h_i}\) , using Algebra, Chapter VII, § 2, no. 1, Proposition 1.

